\ifdefined\gkver\else\def\gkver{student}\fi
\documentclass[\gkver]{gaokaozhenti}
\begin{document}

\chapter{2000年苏浙吉}
%% paperType: 理综物理部分

\begin{enumerate}
\item
%% number: 15
%% paperName: 2000年苏浙吉高考第 15 题 7 分
%% typeId: 1
%% type: 单选题
%% chapter: 原子和原子核
%% point: 天然放射性
%% method: 无
%% score: 7
%% degree: 750
%% duplicateId: 0
%% body:
下面正确的说法是 \xzanswer{C}

① $\beta$ 射线粒子和电子是两种不同的粒子

② 红外线的波长比 X 射线的波长长

③ $\alpha$ 粒子不同于氦原子核

④ $\gamma$ 射线的贯穿本领比 $\alpha$ 粒子的强

\fourchoices[answer=C]
{①②}
{①③}
{②④}
{①④}

%% answer: C
%% memo:
\memoanswer{本题原卷及材料中未提供详解，待补充。}

\item
%% number: 17
%% paperName: 2000年苏浙吉高考第 17 题 7 分
%% typeId: 1
%% type: 单选题
%% chapter: 原子和原子核
%% point: 玻尔模型
%% method: 无
%% score: 7
%% degree: 750
%% duplicateId: 0
%% body:
氢原子的基态能量为 $E_{1}$。下列四个能级图，正确代表氢原子能级的是 \xzanswer{C}

\fourchoices[answer=C, ispicture=true, h=2.8cm]
{figs/02a.png}
{figs/02b.png}
{figs/02c.png}
{figs/02d.png}

%% answer: C
%% memo:
\memoanswer{本题原卷及材料中未提供详解，待补充。}

\item
%% number: 18
%% paperName: 2000年苏浙吉高考第 18 题 7 分
%% typeId: 1
%% type: 单选题
%% chapter: 气体性质
%% point: 液柱移动定性分析
%% method: 无
%% score: 7
%% degree: 550
%% duplicateId: 0
%% body:
一绝热隔板将一绝热长方型容器隔成两部分，两边分别充满气体，隔板可无摩擦移动。开始时，左边的温度为 $0\UdC$，右边的温度为 $20\UdC$，隔板处于静止状态；当左边的气体加热到 $20\UdC$，右边的气体加热到 $40\UdC$ 时，则达到平衡状态时隔板的最终位置 \xzanswer{B}

\fourchoices[answer=B]
{保持不动}
{在初始位置右侧}
{在初始位置左侧}
{决定于加热过程}

%% answer: B
%% memo:
\memoanswer{本题原卷及材料中未提供详解，待补充。}

\item
%% number: 21
%% paperName: 2000年苏浙吉高考第 21 题 7 分
%% typeId: 1
%% type: 单选题
%% chapter: 机械波
%% point: 波的图象
%% method: 无
%% score: 7
%% degree: 550
%% duplicateId: 0
%% body:
如图~\ref{2000苏浙吉04} 所示，图中实线表示横波甲和横波乙在 $t$ 时刻的波形图线，经过 $1\Us$ 后，甲的波峰 $A$ 移到 $A'$ 点，乙的波峰 $B$ 移到 $B'$ 点，如两图中虚线所示。下列说法中正确的是 \xzanswer{D}

① 波甲的波长大于波乙的波长

② 波甲的速度大于波乙的速度

③ 波甲的周期等于波乙的周期

④ 波甲的频率小于波乙的频率

\onepicture[w=8cm, num=4, label=2000苏浙吉04]{figs/04.png}

\fourchoices[answer=D]
{①②}
{②④}
{①④}
{①③}

%% answer: D
%% memo:
\memoanswer{本题原卷及材料中未提供详解，待补充。}

\item
%% number: 29
%% paperName: 2000年苏浙吉高考第 29 题 16 分
%% typeId: 3
%% type: 填空题
%% chapter: 磁场
%% point: 洛伦兹力
%% method: 无
%% score: 16
%% degree: 450
%% duplicateId: 0
%% body:
如图~\ref{2000苏浙吉05} 所示。厚度为 $h$、宽度为 $d$ 的导体板放在垂直于它的磁感应强度为 $B$ 的均匀磁场中，当电流通过导体板时，在导体板的上侧面 $A$ 和下侧面 $A'$ 之间会产生电势差。这种现象称为霍尔效应。实验表明，当磁场不太强时，电势差 $U$、电流 $I$ 和 $B$ 的关系为 $U = k\frac{IB}{d}$，式中的比例系数 $k$ 称为霍尔系数。

霍尔效应可解释如下：外部磁场的洛仑兹力使运动的电子聚集在导体板的一侧，在导体板的另一侧会出现多余的正电荷，从而形成横向电场。横向电场对电子施加与洛仑兹力方向相反的静电力。当静电力与洛仑兹力达到平衡时，导体板上下两侧之间就会形成稳定的电势差。设电流 $I$ 是由电子的定向流动形成的，电子的平均定向速度为 $v$，电量为 $e$。回答下列问题：

\begin{enumerate}
\item
达到稳定状态时，导体板上侧面 $A$ 的电势 \tkanswer{低于} 下侧面 $A'$ 的电势（填“高于”“低于”或“等于”）。
\item
电子所受的洛仑兹力的大小为 \tkanswer{$evB$}。
\item
当导体板上下两侧之间的电势差为 $U$ 时，电子所受静电力的大小为 \tkanswer{$\frac{eU}{h}$}。
\item
由静电力和洛仑兹力平衡的条件，证明霍尔系数为 $k = \frac{1}{ne}$。其中 $n$ 代表导体板单位体积中电子的个数。
\end{enumerate}

\onepicture[w=9cm, num=5, label=2000苏浙吉05, align=right]{figs/05.png}

%% answer: （1）低于；（2）$evB$；（3）$\frac{eU}{h}$（或 $evB$）；（4）证明见解析
\jdanswer{
\begin{enumerate}
\item
低于
\item
$evB$
\item
$\frac{eU}{h}$（或 $evB$）
\item
电子受到横向静电力与洛仑兹力的作用，两力平衡，有 $\frac{eU}{h} = evB$

得 $U = hvB$

通过导体的电流密度 $I = nev\cdot dh$

由 $U = k\frac{IB}{d}$，有 $hvB = k\frac{nevdh\cdot B}{d}$

得 $k = \frac{1}{ne}$
\end{enumerate}
}

%% memo:
\memoanswer{本题原卷及材料中未提供详解，待补充。}

\end{enumerate}

\end{document}
